\documentclass[a4paper,11pt]{article}

\usepackage[utf8]{inputenc}
\usepackage[T1]{fontenc}
\usepackage[german]{babel}
\usepackage[a4paper,margin=2.5cm]{geometry}
\usepackage{enumitem}
\usepackage{titlesec}
\usepackage{longtable}
\usepackage{setspace}
\usepackage{parskip}

\begin{document}

% ---------------------------
% Deckblatt
% ---------------------------

\begin{titlepage}
    \centering

    \vspace*{3cm}

    {\Huge \textbf{Meetingprotokoll}}\\[0.5cm]
    {\Large Statusbericht zur Projektarbeit}\\[2cm]

    \rule{12cm}{0.5pt}\\[0.5cm]

    {\Large Projektarbeit}\\[0.3cm]
    {\large 17. Juli 2026}

    \rule{12cm}{0.5pt}

    \vfill

    {\large THWS / MINOVA}

\end{titlepage}

% ---------------------------
% Inhaltsverzeichnis
% ---------------------------

\tableofcontents
\newpage

% ---------------------------
% Inhalt
% ---------------------------

\section{Allgemeine Informationen}

\begin{tabular}{ll}
\textbf{Datum:} & 17.07.2026 \\
\textbf{Art des Meetings:} & Projektmeeting / Dashboard-Konzeptbesprechung \\
\textbf{Nächstes technisches Meeting:} & 20.07.2026 um 09:00 Uhr \\
\textbf{Nächstes Projektmeeting:} & 06.08.2026 um 09:00 Uhr \\
\end{tabular}

\vspace{1cm}

\section{Projektstatus}

Im Rahmen des Meetings wurde das Konzept des geplanten Avisierungsdashboards gemeinsam besprochen. Der Schwerpunkt lag auf der Überprüfung des Dashboard-Dokuments sowie der zugehörigen Architekturdiagramme.

Das Dokument zum geplanten Avisierungsdashboard wurde insgesamt positiv bewertet. Es wurde bestätigt, dass sich die Ausarbeitung in die richtige Richtung entwickelt. Gleichzeitig wurden offene Punkte sowie Verbesserungsmöglichkeiten hinsichtlich der Diagramme, der Benennung der Systemkomponenten und der Dokumentationsstruktur diskutiert.

\vspace{0.5cm}

\section{Besprochene Punkte}

Während des Meetings wurden folgende Themen behandelt:

\begin{itemize}[leftmargin=1.2cm]

\item Das Dokument zum geplanten Avisierungsdashboard wurde gemeinsam durchgesehen. Unklarheiten sowie erforderliche Anpassungen der Architekturdiagramme wurden gemeinsam besprochen.

\item Für den Abschnitt \textbf{Anforderungserhebung} im Projektbericht soll erläutert werden, welche Anforderungen dem geplanten Avisierungsdashboard zugrunde liegen und welchen Zweck das Dashboard innerhalb des Gesamtsystems erfüllt.

\item Die Architekturdiagramme sollen übersichtlicher und verständlicher gestaltet werden.

\item Die Benennung aller Komponenten soll innerhalb der gesamten Dokumentation konsistent erfolgen. Neue Bezeichnungen sollen eindeutig definiert und anschließend durchgängig verwendet werden.

\item Für die Systemarchitektur wurde folgende Zuordnung festgelegt:

\begin{itemize}
    \item Avisierungsservice $\rightarrow$ Backend
    \item Avisierungsdashboard $\rightarrow$ Frontend
\end{itemize}

\item Sämtliche Diagramme und Abbildungen sollen mit einer Abbildungsnummer (z.\,B. \textit{Abbildung 1}, \textit{Abbildung 2}) sowie einer aussagekräftigen Abbildungsbeschriftung versehen werden.

\item Im Fließtext soll auf die jeweiligen Abbildungen verwiesen werden. Die Bestandteile der Diagramme sollen erläutert werden, sodass deren Funktion innerhalb der Systemarchitektur nachvollziehbar beschrieben wird.

\item Es wurde besprochen, mögliche kostenlose Lösungen für die Integration weiterer Kommunikationskanäle (z.\,B. SMS und WhatsApp) zu recherchieren und hinsichtlich ihrer Eignung für das Projekt zu bewerten.

\end{itemize}

\vspace{0.5cm}

\section{Nächste Schritte}

\begin{itemize}[leftmargin=1.2cm]

\item Überarbeitung des Dokuments zum geplanten Avisierungsdashboard.

\item Überarbeitung der Gliederung der Projektdokumentation.

\item Verbesserung der Übersichtlichkeit und Lesbarkeit der Architekturdiagramme.

\item Recherche möglicher kostenloser Lösungen für die Integration weiterer Kommunikationskanäle (z.\,B. SMS und WhatsApp).

\end{itemize}

\vspace{0.5cm}

\section{Nächste Projektphase}

Die nächste Projektphase konzentriert sich auf die Fertigstellung eines funktionsfähigen Avisierungsdashboards sowie auf die Weiterentwicklung der Backend-Architektur. Parallel dazu sollen mögliche kostenlose Lösungen für die Integration weiterer Kommunikationskanäle recherchiert und bewertet werden.
\vspace{0.5cm}

\section{Nächste Meilensteine}

\begin{itemize}[leftmargin=1.2cm]

\item Fertigstellung eines funktionsfähigen Avisierungsdashboards.

\item Recherche und Bewertung kostenloser Möglichkeiten zur Integration weiterer Kommunikationskanäle (z.\,B. SMS oder WhatsApp).

\item Überarbeitung der Backend-Architektur bis zum technischen Meeting mit Herrn Professor Braun am 20.07.2026.

\end{itemize}

\vspace{1cm}

\section{Zusammenfassung}

Im Meeting wurde das Konzept des geplanten Avisierungsdashboards gemeinsam besprochen und insgesamt positiv bewertet. Es wurde bestätigt, dass sich die Ausarbeitung in die richtige Richtung entwickelt. Gleichzeitig wurden verschiedene Verbesserungen hinsichtlich der Architekturdiagramme, der konsistenten Benennung der Systemkomponenten sowie der Dokumentationsstruktur vereinbart.

Für die nächste Projektphase wurde festgelegt, das Avisierungsdashboard in einen funktionsfähigen Zustand zu überführen, die Backend-Architektur bis zum technischen Review mit Herrn Professor Braun zu überarbeiten und mögliche kostenlose Lösungen für die Integration weiterer Kommunikationskanäle zu recherchieren.

\end{document}